\documentclass[11pt]{article}
\usepackage[margin=1in]{geometry}
\usepackage{amsmath,amssymb,amsfonts}
\usepackage{algorithm}
\usepackage{algpseudocode}
\usepackage[T1]{fontenc}
\usepackage[utf8]{inputenc}
\title{Algorithmic Pseudocode Sample 28}
\author{Source-backed Image2Struct algorithm sample}
\date{}
\begin{document}
\maketitle
\section{Algorithm}
This sample contains algorithmic pseudocode extracted from a source-backed LaTeX benchmark dataset. It is wrapped in a minimal article document for pdfLaTeX validation.

\begin{algorithm}[htbp]
\caption{Source-backed algorithmic procedure}
\begin{algorithmic}
\Function{Transport}{$\bold{M}, \boldsymbol{\beta},\bold{g}, \Delta t$} \Comment{Input moments, natural parameters, and gauge parameters for all cells}
\State $\boldsymbol{\beta}\gets \Call{NewtonOptimization}{\bold{M},\boldsymbol{\beta}, \boldsymbol{\phi}(\bold{u} ;\mathbf{g})}$ \Comment{Solve the natural parameters by Alg. \ref{Newton method} in the gauge $\mathbf{g}$}
\State $\bold{F} \gets \Call{ComputeFluxes}{\boldsymbol{\beta}, \boldsymbol{\phi}(\bold{u} ;\mathbf{g})}$\Comment{Compute fluxes of the statistics $\boldsymbol{\phi}(\bold{u} ;\mathbf{g})$ by \eqref{The ME equation:subeq3} }
\State $\bold{M}_{\pm1}, \bold{F}_{\pm1} \gets \Call{SpatialGaugeTransformation}{\bold{M}, \bold{F}, \bold{g}}$\Comment{Perform gauge transformation as in \eqref{The lax F convert} }
\State $\bold{M} \gets \Call{FiniteVolumeStep}{\bold{M}, \bold{F},\bold{M}_{\pm1}, \bold{F}_{\pm1}, \bold{g}, \Delta t}$\Comment{One step forward in time by \eqref{The ME equation fvm 2} }
\State \textbf{return} $\bold{M}, \boldsymbol{\beta}, \bold{g}$
\EndFunction
\Function{Collision}{$\bold{M}, \boldsymbol{\beta},\bold{g}, \Delta t$} \Comment{Input moments, natural parameters, and gauge parameters for all cells}
\State $\bold{M} \gets \Call{SourceTerm}{\bold{M}, \bold{g}, \Delta t}$\Comment{Compute source term by operator splitting \eqref{source update} }
\State \textbf{return} $\bold{M}, \boldsymbol{\beta}, \bold{g}$
\EndFunction
\Function{Step}{$\bold{M}, \boldsymbol{\beta},\bold{g}$} \Comment{Input moments, natural parameters, and gauge parameters for all cells}
\State $\bold{M}, \boldsymbol{\beta}, \bold{g} \gets \Call{Collision}{\bold{M},\boldsymbol{\beta},\bold{g}, \Delta t/2}$\Comment{Compute collision term by operator splitting}
\State $\bold{M}, \boldsymbol{\beta}, \bold{g} \gets \Call{Transport}{\bold{M},\boldsymbol{\beta},\bold{g}, \Delta t}$\Comment{Compute Transport term by operator splitting}
\State $\bold{M}, \boldsymbol{\beta}, \bold{g} \gets \Call{Collision}{\bold{M},\boldsymbol{\beta},\bold{g}, \Delta t/2}$\Comment{Compute collision term by operator splitting}
\State $\bold{g}_H \gets \Call{ComputeGaugeParameters}{\bold{M}, \bold{g}}$\Comment{Compute the Hermite gauge parameters by \eqref{Hermite parameters in moments} }
\State $\bold{M}_H,\boldsymbol{\beta}_H \gets \Call{GaugeTransformation}{\bold{M}, \boldsymbol{\beta}, \bold{g}_H, \bold{g}}$\Comment{Transform into the Hermite gauge using \eqref{Change gauge M F} and \eqref{Maximal Likelihood equation for Exponential Family Distributions: form transformation}}
\State \textbf{return} $\bold{M}_H, \boldsymbol{\beta}_H, \bold{g}_H$
\EndFunction
\end{algorithmic}
\end{algorithm}
\end{document}
